\documentclass[12pt,letterpaper]{article}
\usepackage[utf8]{inputenc}
\usepackage[spanish]{babel}
\usepackage[margin=2.54cm]{geometry}
\usepackage{mathptmx}
\usepackage{setspace}
\usepackage{enumitem}
\usepackage{hyperref}

\setlength{\parindent}{0pt}
\setlength{\parskip}{6pt}
\singlespacing

\begin{document}

\begin{center}
\textbf{\Large Encuestas que Sí Sirven}
\end{center}

\begin{center}
\textit{P. L. Alvarez Gonzalez} \\
\textit{7690-22-5839} \\
Universidad Mariano Gálvez \\
Seminario de Tecnologías de Información \\
\textit{palvarezg1@miumg.edu.gt}
\end{center}

\vspace{0.5cm}

\textbf{Resumen}

El presente artículo tiene como objetivo explicar, de forma fácil y práctica, cómo se debe diseñar y aplicar una encuesta para que realmente sirva y no se convierta en una pérdida de tiempo tanto para quien la hace como para quien la responde. Para lograrlo se investigó en paginas especializada en metodología de la investigación y se ordenó la información en pasos que cualquier persona puede seguir: definir el objetivo, elegir bien a quién se le va a preguntar, hacer las preguntas claras, escoger el tipo adecuado y hacer una prueba piloto antes de lanzar la versión final. También se toman los errores más comunes que la gente comete al armar una encuesta, como preguntas ambiguas, preguntas que inducen la respuesta o cuestionarios demasiado largos. Como resultado se logró construir una guía práctica que sirve tanto para trabajos académicos como para encuestas reales, por ejemplo en plataformas digitales como Google Forms. Se concluye que una encuesta efectiva no depende de la cantidad de preguntas que tenga, sino de qué tan bien fueron pensadas antes de escribirlas.

\textbf{Palabras clave:} encuesta, cuestionario, muestra, metodología, recolección de datos.

\textbf{Desarrollo del tema}

Una encuesta es, en pocas palabras, una forma de preguntarle a un grupo de personas sobre lo que piensan, saben o hacen respecto a un tema. Dentro de los métodos de investigación, la encuesta se apoya en dos herramientas principales: la entrevista, que es una conversación directa entre el investigador y la persona, y el cuestionario, que es un formulario que la persona llena por su cuenta. El cuestionario es el más usado hoy en día porque es barato, se puede aplicar a mucha gente en poco tiempo y es fácil de tabular, sobre todo con herramientas digitales.

El primer paso para hacer una encuesta efectiva es tener bien claro el objetivo. Antes de escribir una sola pregunta hay que responder: ¿qué es lo que se quiere saber? Si el objetivo no está claro, el cuestionario termina siendo una mezcla de preguntas sueltas que no llevan a ninguna conclusión útil. Por ejemplo, no es lo mismo querer medir el conocimiento de un tema que querer medir la opinión sobre ese mismo tema; ambas cosas requieren preguntas distintas.

El segundo paso es definir bien a quién se le va a preguntar, es decir, la población y la muestra. No tiene sentido encuestar a cualquiera si el estudio busca información específica de un grupo en particular. La muestra debe ser representativa, o sea, debe reflejar las características del grupo completo que se quiere estudiar, aunque no se pueda encuestar a todos.

El tercer paso, y probablemente el más importante, es hacer  las preguntas. Aquí es donde más se equivocan las personas que hacen encuestas por primera vez. Las preguntas deben ser cortas, claras y deben preguntar una sola cosa a la vez. Un error muy común es hacer preguntas dobles, por ejemplo: "¿Le gusta la comida rápida y la considera saludable?", porque la persona puede estar de acuerdo con una parte y no con la otra, y no sabe qué responder. También hay que evitar preguntas que ya traen la respuesta incluida, como ¿No cree usted que...?, porque eso influye en lo que la persona va a contestar.

El tipo de pregunta también importa. Existen preguntas cerradas, donde la persona elige entre opciones ya definidas, y preguntas abiertas, donde puede responder con sus propias palabras. Las cerradas son más fáciles de analizar, mientras que las abiertas dan información más rica pero cuestan más de procesar. Dentro de las cerradas, las de rango tipo (por ejemplo, del 1 al 5, donde 1 es totalmente en desacuerdo y 5 es totalmente de acuerdo son muy usadas porque permiten medir opiniones o actitudes de forma numérica.

El orden de las preguntas también influye en los resultados. Se recomienda iniciar con preguntas generales o fáciles de responder, dejar las preguntas más sensibles o personales para el final, y agrupar las preguntas por temas para que la persona no se pierda ni se aburra a la mitad del formulario.

Antes de lanzar la encuesta de forma masiva, es clave hacer una prueba piloto, es decir, aplicar el cuestionario a un grupo pequeño para detectar preguntas confusas, errores de redacción o preguntas que nadie entiende. Esto evita tener que corregir la encuesta después de haber recolectado muchas respuestas, lo cual sería una pérdida de tiempo enorme.

Finalmente, que tan grande es el cuestionario también es un factor que afecta la efectividad. Entre más larga sea la encuesta, más alta es la probabilidad de que la persona la abandone a la mitad o conteste sin pensar solo para terminar rápido. Por eso se recomienda ser lo más breve posible sin sacrificar la información que realmente se necesita.

\textbf{Observaciones y comentarios}

Durante esta investigación se pudo notar que la mayoría de errores en las encuestas no vienen de la falta de conocimiento técnico, sino de la prisa. Muchas personas escriben las preguntas pensando en lo que a ellos les parece obvio, sin ponerse en el lugar de quien va a responder, y eso genera confusión. También se observó que las herramientas digitales como Google Forms facilitan mucho el proceso de aplicación , pero no garantizan que la encuesta esté bien diseñada. La tecnología ayuda con la logística, no con el contenido. 

\textbf{Conclusiones}

\begin{enumerate}[label=\arabic*.]
\item Una encuesta efectiva empieza con un objetivo claro, sin eso, ninguna pregunta bien redactada sirve de mucho.
\item La forma en que se redactan las preguntas influye directamente en la calidad de las respuestas, por lo que se debe evitar la ambigüedad y las preguntas que inducen una respuesta.
\item Hacer una prueba piloto antes de aplicar la encuesta de forma masiva ahorra tiempo y evita errores costosos.
\item Las herramientas digitales facilitan la aplicación y el análisis, pero no sustituyen un buen diseño metodológico previo.
\end{enumerate}

\textbf{Referencias}

Hernández-Sampieri, R., Fernández-Collado, C., \& Baptista-Lucio, P. (2014). \textit{Metodología de la investigación} (6.ª ed.). McGraw-Hill.

Malhotra, N. K. (2008). \textit{Investigación de mercados} (5.ª ed.). Pearson Educación.

American Psychological Association. (2020). \textit{Publication manual of the American Psychological Association} (7.ª ed.).

\vspace{0.5cm}
\noindent \small Repositorio: https://github.com/PLAlvarez11/seminario-encuesta

\noindent \small Encuesta: https://forms.gle/yncD3SzuhsL237Ra9


\end{document}